% figures/f2_deletion.tex -- the two deletion cases worked on graphs (k = 1).
% Rows of R are vertices with keys 0 < 1 < 2 < 3 (top to bottom, so the low of
% a column is its lowest non-zero entry); rows of V are edges; columns are slots.
% Every entry is computed and checked by reference/figure_examples.py.
\newcommand{\pv}[1]{{\setlength{\fboxsep}{0.6pt}\setlength{\fboxrule}{0.6pt}\fbox{$#1$}}}
\newcommand{\fc}{\cellcolor{black!14}}
\newcommand{\chg}[1]{\boldsymbol{#1}}
\tikzset{
  dvtx/.style={circle, draw=black!85, fill=white, inner sep=0pt, minimum size=3.6mm,
    font=\scriptsize},
  dedg/.style={black!80, semithick},
  dsig/.style={black, very thick, dash pattern=on 2.4pt off 1.5pt},
}
\newcommand{\whynote}[1]{{\scriptsize\raggedright #1\par}}
{\footnotesize\renewcommand{\arraystretch}{1.06}%
%
% =========================== (a) Case A ===================================
\begin{minipage}[t]{0.485\textwidth}
\textbf{(a) Case A:} $\sigma=01$ lies on a cycle\par\smallskip
\begin{tabular}{@{}c@{\hspace{0.7em}}l@{\hspace{1.1em}}l@{}}
\begin{tikzpicture}[baseline=(current bounding box.center)]
  \node[dvtx] (a0) at (0,0) {0}; \node[dvtx] (a1) at (1.1,0) {1};
  \node[dvtx] (a2) at (0.55,0.95) {2};
  \draw[dedg] (a0) -- (a2) (a1) -- (a2);
  \draw[dsig] (a0) -- node[below=-1pt, font=\scriptsize] {$\sigma$} (a1);
  \node[font=\scriptsize] at (0.55,-0.62) {$\beta=(1,1)$};
\end{tikzpicture}
& $V=\begin{array}{@{}r|ccc@{}}
    & s_1 & s_2 & s_3\\ \hline
    \fc 01 & \fc 1 & \fc 0 & \fc 1\\
    02 & 0 & 1 & 1\\
    12 & 0 & 0 & 1\end{array}$
& $R=\begin{array}{@{}r|ccc@{}}
    & s_1 & s_2 & s_3\\ \hline
    0 & 1 & 1 & 0\\
    1 & \pv{1} & 0 & 0\\
    2 & 0 & \pv{1} & 0\end{array}$\\[16pt]
$\Big\downarrow$ & \multicolumn{2}{@{}l@{}}{\parbox[c]{5.75cm}{\scriptsize\raggedright
    $O=\occ(\sigma)=\{s_1,s_3\}$ (the $1$s in the shaded row); $R_{s_3}=0$, so
    $Z=\{s_3\}$ and $z=s_3$.\\
    $V_{s_1}\gets V_{s_1}+V_{s_3}$; remove slot $s_3$ and row $\sigma$.}}\\[14pt]
\begin{tikzpicture}[baseline=(current bounding box.center)]
  \node[dvtx] (b0) at (0,0) {0}; \node[dvtx] (b1) at (1.1,0) {1};
  \node[dvtx] (b2) at (0.55,0.95) {2};
  \draw[dedg] (b0) -- (b2) (b1) -- (b2);
  \node[font=\scriptsize] at (0.55,-0.62) {$\beta=(1,0)$};
\end{tikzpicture}
& $V=\begin{array}{@{}r|cc@{}}
    & s_1 & s_2\\ \hline
    02 & \chg{1} & 1\\
    12 & \chg{1} & 0\end{array}$
& $R=\begin{array}{@{}r|cc@{}}
    & s_1 & s_2\\ \hline
    0 & 1 & 1\\
    1 & \pv{1} & 0\\
    2 & 0 & \pv{1}\end{array}$
\end{tabular}\par\medskip
\whynote{\textbf{Why a cycle.} $\bd V_{s_3}=R_{s_3}=0$, so the addition removes
  $\sigma$ from $V_{s_1}$ and leaves every $R_j$, hence every low, as it was.
  $n_1$ drops and $r_1=2$ stays: $\beta_1$ falls by one.}
\end{minipage}\hfill
%
% =========================== (b) Case B ===================================
\begin{minipage}[t]{0.485\textwidth}
\textbf{(b) Case B:} $\sigma=03$ lies on no cycle\par\smallskip
\begin{tabular}{@{}c@{\hspace{0.7em}}l@{\hspace{1.1em}}l@{}}
\begin{tikzpicture}[baseline=(current bounding box.center)]
  \node[dvtx] (c3) at (0.6,1.0) {3};
  \node[dvtx] (c0) at (0,0) {0}; \node[dvtx] (c1) at (0.6,0) {1};
  \node[dvtx] (c2) at (1.2,0) {2};
  \draw[dedg] (c1) -- (c3) (c2) -- (c3);
  \draw[dsig] (c0) -- node[left=-1pt, font=\scriptsize] {$\sigma$} (c3);
  \node[font=\scriptsize] at (0.6,-0.62) {$\beta=(1,0)$};
\end{tikzpicture}
& $V=\begin{array}{@{}r|ccc@{}}
    & s_1 & s_2 & s_3\\ \hline
    \fc 03 & \fc 1 & \fc 1 & \fc 1\\
    13 & 0 & 1 & 0\\
    23 & 0 & 0 & 1\end{array}$
& $R=\begin{array}{@{}r|ccc@{}}
    & s_1 & s_2 & s_3\\ \hline
    0 & 1 & 1 & 1\\
    1 & 0 & \pv{1} & 0\\
    2 & 0 & 0 & \pv{1}\\
    3 & \pv{1} & 0 & 0\end{array}$\\[16pt]
$\Big\downarrow$ & \multicolumn{2}{@{}l@{}}{\parbox[c]{5.75cm}{\scriptsize\raggedright
    $O=\{s_1,s_2,s_3\}$ and every $R_j\ne0$, so $Z=\emptyset$; the lows are
    $3,1,2$, and $p=s_2$ has the smallest.\\
    $(R_j,V_j)\gets(R_j+R_{s_2},V_j+V_{s_2})$ for $j=s_1,s_3$; remove slot $s_2$
    and row $\sigma$.}}\\[14pt]
\begin{tikzpicture}[baseline=(current bounding box.center)]
  \node[dvtx] (d3) at (0.6,1.0) {3};
  \node[dvtx] (d0) at (0,0) {0}; \node[dvtx] (d1) at (0.6,0) {1};
  \node[dvtx] (d2) at (1.2,0) {2};
  \draw[dedg] (d1) -- (d3) (d2) -- (d3);
  \node[font=\scriptsize] at (0.6,-0.62) {$\beta=(2,0)$};
\end{tikzpicture}
& $V=\begin{array}{@{}r|cc@{}}
    & s_1 & s_3\\ \hline
    13 & \chg{1} & \chg{1}\\
    23 & 0 & 1\end{array}$
& $R=\begin{array}{@{}r|cc@{}}
    & s_1 & s_3\\ \hline
    0 & \chg{0} & \chg{0}\\
    1 & \chg{1} & \chg{1}\\
    2 & 0 & \pv{1}\\
    3 & \pv{1} & 0\end{array}$
\end{tabular}\par\medskip
\whynote{\textbf{Why the smallest low.} Adding $R_{s_1}=\{0,3\}$, whose low $3$
  is the largest, would instead give $\{1,3\}$ and $\{2,3\}$: two lows $3$,
  violating (I2). $R_{s_2}$ lies entirely above the other lows, so adding it
  moves none. $n_1$ and $r_1$ both drop: $\beta_0$ rises by one.}
\end{minipage}}%
